% Part: incompleteness
% Chapter: incompleteness-provability
% Section: fixed-point-lemma

\documentclass[../../../include/open-logic-section]{subfiles}

\begin{document}

\olfileid{inc}{inp}{fix}

\olsection{স্থির-বিন্দু লেমা}

\begin{explain}
স্থির-বিন্দু লেমা বলে, যেকোনো !!{formula}~$!B(x)$-এর জন্য এমন
!!a{sentence}~$!A$ আছে যাতে $\Th{T} \Proves !A \liff !B(\gn{!A})$,
যদি $\Th{T}$ তত্ত্বটি $\Th{Q}$-কে প্রসারিত করে। মিথ্যাবাদী বাক্যের
ক্ষেত্রে আমরা চাইব $!A$ যেন~``$\gn{!A}$ মিথ্যা''-র সমতুল্য হয়
($\Th{T}$-তে প্রমাণযোগ্যভাবে), অর্থাৎ $\Gn{!A}$ কোনো মিথ্যা
!!{sentence}-এর গ্যোডেল সংখ্যা—এই উক্তিটির সমতুল্য। প্রমাণের ধারণাটি
বোঝার জন্য $!A$-সম্বন্ধে কুয়াইনের অনানুষ্ঠানিক ব্যাখ্যার সঙ্গে তুলনা করা
উপকারী: ``{}`নিজের উদ্ধৃত রূপটি আগে বসালে মিথ্যা দেয়'—এর আগে তার নিজের
উদ্ধৃত রূপটি বসালে মিথ্যা দেয়।'' কোনো রাশি নিয়ে তার আগে সেই রাশিরই
উদ্ধৃত রূপ বসিয়ে একটি বাক্য গঠনের প্রক্রিয়াকে ওই রাশিটির
\emph{কর্ণীকরণ} এবং ফলটিকে তার কর্ণীকৃত রূপ বলা যেতে পারে। তাই
`নিজের উদ্ধৃত রূপটি আগে বসালে মিথ্যা দেয়'-এর কর্ণীকৃত রূপ হলো
``{}`নিজের উদ্ধৃত রূপটি আগে বসালে মিথ্যা দেয়'—এর আগে তার নিজের উদ্ধৃত
রূপটি বসালে মিথ্যা দেয়।'' এবার খেয়াল করো, কুয়াইনের মিথ্যাবাদী বাক্যটি
`মিথ্যা দেয়'-এর কর্ণীকৃত রূপ নয়; সেটি `নিজের উদ্ধৃত রূপটি আগে বসালে
মিথ্যা দেয়'-এর কর্ণীকৃত রূপ। অতএব যে ধর্মটির কর্ণীকরণে মিথ্যাবাদী বাক্যটি
পাওয়া যায়, সেই ধর্মের মধ্যেই কর্ণীকরণ জড়িত!{}

পাটিগণিতের ভাষায় একটি মুক্ত চল-সহ !!a{formula}-র গ্যোডেল সংখ্যা গণনা
করে এবং সেই গ্যোডেল সংখ্যার প্রমিত সংখ্যাপদটি মুক্ত চলের জায়গায় বসিয়ে
তার উদ্ধৃত রূপ গড়ি। $!E(x)$-এর কর্ণীকৃত রূপ হলো $!E(\num{n})$, যেখানে
$n = \Gn{!E(x)}$। (এখন থেকে $\num{\Gn{!E(x)}}$-কে সংক্ষেপে
$\gn{!E(x)}$ লিখব।) তাই $!B(x)$ যদি ``মিথ্যা'' হয়, তবে ``নিজের উদ্ধৃত
রূপটি আগে বসালে মিথ্যা দেয়'' হবে ``তার কর্ণীকৃত রূপের গ্যোডেল সংখ্যায়
প্রয়োগ করলে মিথ্যা দেয়।'' $\fn{diag}(n)$ অপেক্ষকটির জন্য যদি
$\Obj{diag}$ প্রতীক থাকত—যে অপেক্ষকটি গ্যোডেল সংখ্যা~$n$-বিশিষ্ট
!!{formula}-টির কর্ণীকৃত রূপের গ্যোডেল সংখ্যা গণনা করে—তবে $!E(x)$-কে
$!B(\Obj{diag}(x))$ লেখা যেত। তখন কুয়াইনের মিথ্যাবাদী বাক্যের রূপটি হতো
তার কর্ণীকরণ, অর্থাৎ $!E(\gn{!E(x)})$ বা
$!B(\Obj{diag}(\gn{!B(\Obj{diag}(x))}))$। অবশ্য এখন $!B(x)$ অন্য যেকোনো
ধর্ম হতে পারে, এবং একই নির্মাণ কাজ করবে। অসম্পূর্ণতা উপপাদ্যের জন্য
$!B(x)$ হিসেবে নেব ``$x$~$\Th{T}$-তে !!{derivable} নয়।'' তখন $!E(x)$
হবে ``তার কর্ণীকৃত রূপের গ্যোডেল সংখ্যায় প্রয়োগ করলে এমন
!!a{sentence} দেয় যা $\Th{T}$-তে !!{derivable} নয়।''

$\Th{T}$-তে এটি বিধিবদ্ধ করতে $\fn{diag}$-কে বিধিবদ্ধ করার উপায় খুঁজতে
হবে। $\fn{diag}(n)$ অপেক্ষকটি গণনসাধ্য, এমনকি আদিম পুনরাবৃত্ত: $n$ যদি
$!E(x)$ সূত্রের গ্যোডেল সংখ্যা হয়, তবে $\fn{diag}(n)$ ফেরত দেয়
$!E(\gn{!E(x)})$-এর গ্যোডেল সংখ্যা। (মনে রেখো, $\gn{!E(x)}$ হলো
$!E(x)$-এর গ্যোডেল সংখ্যার প্রমিত সংখ্যাপদ, অর্থাৎ
$\num{\Gn{!E(x)}}$।) $\Obj{diag}$ যদি $\fn{diag}$ অপেক্ষককে
উপস্থাপনকারী $\Th{T}$-এর কোনো অপেক্ষক-প্রতীক হতো, তবে $!A$ হিসেবে
$!B(\Obj{diag}(\gn{!B(\Obj{diag}(x))}))$ সূত্রটি নেওয়া যেত। খেয়াল করো,
\begin{align*}
\fn{diag}(\Gn{!B(\Obj{diag}(x))}) & =
\Gn{!B(\Obj{diag}(\gn{!B(\Obj{diag}(x))}))} \\
& = \Gn{!A}.
\end{align*}
ধরি $\Th{T}$ নিচেরটি !!{derive} করতে পারে:
\[
\Obj{diag}(\gn{!B(\Obj{diag}(x))}) = \gn{!A},
\]
তাহলে সেটি $!B(\Obj{diag}(\gn{!B(\Obj{diag}(x))}))
\liff !B(\gn{!A})$-ও !!{derive} করতে পারে। কিন্তু সংজ্ঞা অনুসারে বাঁদিকটি
হলো~$!A$।

অবশ্য $\Obj{diag}$ সাধারণভাবে $\Th{T}$-এর অপেক্ষক-প্রতীক হবে না, এবং
নিশ্চিতভাবেই $\Th{Q}$-এর কোনো প্রতীক নয়। কিন্তু $\fn{diag}$ গণনসাধ্য
বলে কোনো সূত্র $!D_{\fn{diag}}(x,y)$ দ্বারা সেটি $\Th{Q}$-তে
\emph{উপস্থাপনযোগ্য}। তাই $!B(\Obj{diag}(x))$ না লিখে আমরা
$\lexists[y][(!D_{\fn{diag}}(x,y) \land !B(y))]$ লিখতে পারি। তা ছাড়া
উপরের প্রমাণের রূপরেখাটি অপরিবর্তিত থাকে; বস্তুত, সেটি $\Th{Q}$-তেই
কার্যকর।
\end{explain}

\begin{lem}
\ollabel{lem:fixed-point} ধরি $!B(x)$ হলো কেবল $x$ মুক্ত এমন যেকোনো
সূত্র। তবে এমন একটি বাক্য~$!A$ আছে যাতে $\Th{Q} \Proves
!A \liff !B(\gn{!A})$।
\end{lem}


\begin{proof}
$!B(x)$ দেওয়া থাকলে $!E(x)$ হিসেবে
$\lexists[y][(!D_{\fn{diag}}(x,y) \land !B(y))]$ সূত্রটি নিই এবং $!A$-কে
তার কর্ণীকৃত রূপ, অর্থাৎ $!E(\gn{!E(x)})$ সূত্রটি ধরি।

$!D_{\fn{diag}}$ যেহেতু $\fn{diag}$-কে উপস্থাপন করে এবং
$\fn{diag}(\Gn{!E(x)}) = \Gn{!A}$, তাই $\Th{Q}$ নিচেরগুলি !!{derive}
করতে পারে:
\begin{align}
  & !D_{\fn{diag}}(\gn{!E(x)}, \gn{!A}) \ollabel{repdiag1} \\
  & \lforall[y][(!D_{\fn{diag}}(\gn{!E(x)},y) \lif
  \eq[y][\gn{!A}])]. \ollabel{repdiag2}
\end{align}
এবার দেখাব $\Th{Q} \Proves !A \liff !B(\gn{!A})$। শুধু যুক্তিবিদ্যা ও
$\Th{Q}$-তে !!{derivable} তথ্য ব্যবহার করে অনানুষ্ঠানিকভাবে যুক্তি দিই।

প্রথমে $!A$, অর্থাৎ $!E(\gn{!E(x)})$, ধরে নিই। $!E(x)$-এর সংজ্ঞায়
ফিরলে দেখি, $!E(\gn{!E(x)})$ ঠিক নিচের সূত্রটি:
\[
\lexists[y][(!D_{\fn{diag}}(\gn{!E(x)},y) \land !B(y))].
\]
এমন একটি~$y$ বিবেচনা করো। যেহেতু
$!D_{\fn{diag}}(\gn{!E(x)},y)$, তাই \olref{repdiag2} অনুসারে
$y = \gn{!A}$। অতএব $!B(y)$ থেকে পাই~$!B(\gn{!A})$।

এবার $!B(\gn{!A})$ ধরে নিই। \olref{repdiag1} অনুসারে পাই
\begin{align*}
& !D_{\fn{diag}}(\gn{!E(x)}, \gn{!A}) \land !B(\gn{!A}).
\intertext{এর থেকে অনুসৃত হয়}
& \lexists[y][(!D_{\fn{diag}}(\gn{!E(x)},y) \land !B(y))].
\end{align*}
কিন্তু এটি ঠিক $!E(\gn{!E(x)})$, অর্থাৎ~$!A$।
\end{proof}

\begin{digress}
গণনসাধ্যতা তত্ত্বে স্থির-বিন্দু লেমার প্রমাণের সঙ্গে এই প্রমাণটি তুলনা
করো। পার্থক্য হলো, এখানে আমরা নিজেকে ব্যবহার করে একটি \emph{উক্তি}
সংজ্ঞায়িত করতে চাই, আর সেখানে নিজেকে ব্যবহার করে একটি \emph{অপেক্ষক}
সংজ্ঞায়িত করতে চেয়েছিলাম; এই পার্থক্য বাদ দিলে ধারণাটি সত্যিই একই।
\end{digress}

\begin{prob}
  !!^a{formula}~$!A(x)$ একটি \emph{সত্যতার সংজ্ঞা}, যদি প্রতিটি
  !!{sentence}s~$!B$-এর জন্য $\Th{Q} \Proves !B \liff !A(\gn{!B})$।
  স্থির-বিন্দু লেমা ব্যবহার করে দেখাও যে কোনো !!{formula}-ই সত্যতার
  সংজ্ঞা নয়।
\end{prob}

\end{document}
